%% kapa: exam=1 round=35 year=2024 period=1 form=A subject=civil_law subject_ko=민법 seq=21 number=21
%% source: Q-Net 감정평가사 기출문제 (cst00309) artlSeq=5213229 file=raw/1008/5213229/1교시 문제지 원본.pdf
%% conversion: kordoc 4.24.51 -> tools/gen1.py (auto); review=proofread-sonnet; flags=none
\begin{question}{21}
물권에 관한 설명으로 옳지 않은 것은? (다툼이 있으면 판례에 따름)
\begin{choices}
  \item 적법한 분할절차를 거치지 않은 채 토지 중 일부만에 관하여 소유권보존등기를 할 수 없다.
  \item 온천에 관한 권리는 관습법상의 물권이 아니다.
  \item 1필 토지의 일부도 점유취득시효의 대상이 될 수 있다.
  \item 부속건물로 등기된 창고건물은 분할등기 없이 원채인 주택과 분리하여 경매로 매각될 수 있다.
  \item 지상권은 저당권의 객체가 될 수 있다.
\end{choices}
\end{question}
\answer{4}
